% This file provides an example Beamer presentation using the RWTH theme
% showcasing some of the more common options, similar to the Powerpoint version
% 12.11.2014: Revision 1 (Harold Bruintjes, Tim Lange)

% For RWTH, beamer should be loaded with class option t (top)
\documentclass[aspectratio=169]{beamer}
\newcommand{\PROPOSAL}{Proposal}

% Use fontspec to get Arial font
% Requires use of XeLaTeX
%\usepackage{fontspec}
%\setmainfont{Arial}
%\setsansfont{Arial}
% Also force Arial for math for a more consistent look
%\usepackage{unicode-math}
%\setmathfont{Arial}

\usepackage{helvet}
\renewcommand{\familydefault}{\sfdefault}
\usepackage{lmodern}%
\usepackage{MnSymbol}%
\usepackage[utf8]{inputenc} % Allows inpit in utf8
% German style date formatting (footer)
\usepackage[ddmmyyyy]{datetime}

\usepackage{booktabs} % Allows the use of \toprule, \midrule and \bottomrule in tables

\renewcommand{\dateseparator}{.}

% Format the captions used for figures etc.
\usepackage[compatibility=false]{caption}
%\captionsetup{singlelinecheck=off,justification=raggedleft,labelformat=empty,labelsep=none}

% PGFPlots is used for drawing some of the charts
\usepackage{pgfplots}
\pgfplotsset{compat=newest}
% This file contains some styles and macros for drawing charts similar to those of MS Office

\pgfplotsset{hor_barchart/.style={
  xbar=0mm,
  xmin=0,
  xtick=\empty,
  axis y line*=left,
  x axis line style={opacity=0},
  bar width=0.6cm,
  ytick=data,
  nodes near coords,
  every axis/.append style={font=\normalsize},
  every node near coord/.append style={font=\normalsize},
  nodes near coords align={horizontal},
  legend style={at={(0,-10mm)},anchor=north west,legend columns=-1,draw=none},
}}

\pgfplotsset{ver_barchart/.style={
  ybar=0mm,
  x = 4.5cm,
  ymin=0,
  ymajorgrids,
  axis x line*=bottom,
  y axis line style={opacity=0},
  bar width=0.8cm,
  enlarge x limits={0.15},
  xtick=data,
  nodes near coords,
  every axis/.append style={font=\normalsize},
  every node near coord/.append style={font=\normalsize},
  nodes near coords align={vertical},
  legend style={at={(0,-10mm)},anchor=north west,legend columns=-1,draw=none},
}}

\tikzstyle{chart}=[
    legend label/.style={font={\normalsize},anchor=west,align=left},
    legend box/.style={rectangle, draw=none, minimum size=5pt},
]

\tikzstyle{pie chart}=[
    chart,
    slice/.style={line cap=round, line join=round,draw=none},
    pie title/.style={font={\bf}},
    slice type/.style 2 args={
        ##1/.style={fill=##2},
        values of ##1/.style={}
    }
]

\newcommand{\pie}[3][]{
    \begin{scope}[#1]
    \pgfmathsetmacro{\curA}{90}
    \pgfmathsetmacro{\r}{1}
    \def\c{(0,0)}
    \node[pie title] at (90:1.3) {#2};
    \foreach \v/\s in{#3}{
        \pgfmathsetmacro{\deltaA}{\v/100*360}
        \pgfmathsetmacro{\nextA}{\curA + \deltaA}
        \pgfmathsetmacro{\midA}{(\curA+\nextA)/2}

        \path[slice,\s] \c
            -- +(\curA:\r)
            arc (\curA:\nextA:\r)
            -- cycle;

        %\begin{pgfonlayer}{foreground}
        % Position labels at 1.2 times radius (just outside of chart)
        \path \c -- node[pos=1.2,pie values,values of \s]{$\v\%$} +(\midA:\r);
        %\end{pgfonlayer}

        \global\let\curA\nextA
    }
    \end{scope}
}

% Custom legend (used for pie chart)
\newcommand{\legend}[2][]{
\begin{scope}[#1]
  \path
    \foreach \n/\s in {#2} {
      ++(0,-5pt) node[\s,legend box] {} +(5pt,0) node[legend label] {\n}
    };
\end{scope}
}


% Load the actual RWTH theme. Suggested is to load the full theme,
% as it requires some specific dimensions
\usetheme{rwth}
% comment in this line if you are writing your bachelor thesis
%\newcommand*{\BACHELOR}{}
% Adjust your information
\title{Thesis Title}
\subtitle{\ifdefined\BACHELOR{} Bachelor \else Master \fi Thesis \ifdefined\PROPOSAL{}
	\PROPOSAL{}
	\fi}
\date{April 6, 2020}
\newcommand{\firstname}{firstname}
\newcommand{\lastname}{lastname}
\newcommand{\matrNo}{123456}
\newcommand{\email}{your.mail@rwth-aachen.de}
\newcommand{\studyProgram}{\ifdefined\BACHELOR{} Bachelor \else Master \fi Computer Science}

% Adjust first supervisor information
\newcommand{\firstsupervisor}{first supervisor}
\newcommand{\firstsupervisorchair}{chair of first supervisor}
\newcommand{\firstsupervisoruniversity}{RWTH Aachen University}

% Adjust second supervisor information
\newcommand{\secondsupervisor}{second supervisor}
\newcommand{\secondsupervisorchair}{chair of second supervisor}
\newcommand{\secondsupervisoruniversity}{RWTH Aachen University}


% Adjust advisor information
\newcommand{\firstadvisor}{first advisor}
\newcommand{\firstadvisorchair}{chair of first advisor}
\newcommand{\firstadvisoruniversity}{RWTH Aachen University}

% comment these lines out if you have only one advisor
\newcommand{\secondadvisor}{second advisor}
\newcommand{\secondadvisorchair}{chair of second advisor}
\newcommand{\secondadvisoruniversity}{RWTH Aachen University}
\graphicspath{{../../figures/}}

% Setup presentation information
\author[\firstname]{\firstname~\lastname~\email}
\institute[RWTH]{RWTH Aachen University}



% Set the logo to the file `logo`
% It will be scaled automatically
\logo{\includegraphics{rwth_i5_en_rgb}}

% Uncomment this if you want a TOC at every section start
%\AtBeginSection{\frame{
%    \frametitle{Content}
%    \tableofcontents[currentsection]
%}}

% Use this to control some aspects of the footer
\setbeamertemplate{footertextextra}{\ifdefined\BACHELOR Bachelor \else Master \fi Thesis \ifdefined\PROPOSAL
\PROPOSAL
\fi}
%\setbeamertemplate{footertext}{Example of an overridden footer}

\begin{document}

% Note: Title pages should be created as plain
% Title page with a blue bar
\setbeamercolor{title page bar}{fg=rwth}
\setbeamertemplate{title page}[rwth]{}
\begin{frame}[plain]
\titlepage

\parbox{0cm}{
  \vspace{-5.5em}
  \begin{tabbing}
    \=\textbf{\firstname~\lastname}\=\\[0.5em]
    \=Study Program: \studyProgram~~~\=Matr.-Nr.:\matrNo\\[0.5em]
    \=Chair of Computer Science 5 Information Systems \& Databases \=\\[0.5em]
    \={Supervisors}:~~~\=\firstsupervisor\\
    \>                 \>\secondsupervisor\\[0.5em]
    \>{Advisor}:       \>\firstadvisor\ifdefined\secondadvisor\\
    \>                 \>\secondadvisor\\[0.5em]
    \fi
  \end{tabbing}
}
\end{frame}

% Title page with a 1/3rd size picture
\setbeamercolor{title page bar}{fg=white}
\setbeamertemplate{title page}[rwth][title_small]{}
\begin{frame}[plain]
\vspace{1cm}
\titlepage

\vspace{-6em}
\parbox{0cm}{
  \begin{tabbing}
    \=\textbf{\firstname~\lastname}\=\\[0.4em]
    \=Study Program: \studyProgram~~~\=Matr.-Nr.:\matrNo\\[0.4em]
    \=Chair of Computer Science 5 Information Systems \& Databases \=\\[0.4em]
    \={Supervisors}:~~~\=\firstsupervisor\\
    \>                 \>\secondsupervisor\\[0.4em]
    \>{Advisor}:       \>\firstadvisor\ifdefined\secondadvisor\\
    \>                 \>\secondadvisor\\[0.4em]
    \fi
  \end{tabbing}
}
\end{frame}

% Frame with items
\section*{Overview}
\begin{frame}
  \tableofcontents % Throughout your presentation, if you choose to use \section{} and \subsection{} commands, these will automatically be printed on this slide as an overview of your presentation
\end{frame}

\section{First Section} % Sections can be created in order to organize your presentation into discrete blocks, all sections and subsections are automatically printed in the table of contents as an overview of the talk
%------------------------------------------------

\subsection{Subsection Example} % A subsection can be created just before a set of slides with a common theme to further break down your presentation into chunks

\begin{frame}
\frametitle{Paragraphs of Text}
Sed iaculis dapibus gravida. Morbi sed tortor erat, nec interdum arcu. Sed id lorem lectus. Quisque viverra augue id sem ornare non aliquam nibh tristique. Aenean in ligula nisl. Nulla sed tellus ipsum. Donec vestibulum ligula non lorem vulputate fermentum accumsan neque mollis.\\~\\

Sed diam enim, sagittis nec condimentum sit amet, ullamcorper sit amet libero. Aliquam vel dui orci, a porta odio. Nullam id suscipit ipsum. Aenean lobortis commodo sem, ut commodo leo gravida vitae. Pellentesque vehicula ante iaculis arcu pretium rutrum eget sit amet purus. Integer ornare nulla quis neque ultrices lobortis. Vestibulum ultrices tincidunt libero, quis commodo erat ullamcorper id.
\end{frame}

%------------------------------------------------

\begin{frame}
\frametitle{Bullet Points}
\begin{itemize}
\item Lorem ipsum dolor sit amet, consectetur adipiscing elit
\item Aliquam blandit faucibus nisi, sit amet dapibus enim tempus eu
\item Nulla commodo, erat quis gravida posuere, elit lacus lobortis est, quis porttitor odio mauris at libero
\item Nam cursus est eget velit posuere pellentesque
\item Vestibulum faucibus velit a augue condimentum quis convallis nulla gravida
\end{itemize}
\end{frame}

%------------------------------------------------

\begin{frame}
\frametitle{Blocks of Highlighted Text}
\begin{block}{Block 1}
Lorem ipsum dolor sit amet, consectetur adipiscing elit. Integer lectus nisl, ultricies in feugiat rutrum, porttitor sit amet augue. Aliquam ut tortor mauris. Sed volutpat ante purus, quis accumsan dolor.
\end{block}

\begin{block}{Block 2}
Pellentesque sed tellus purus. Class aptent taciti sociosqu ad litora torquent per conubia nostra, per inceptos himenaeos. Vestibulum quis magna at risus dictum tempor eu vitae velit.
\end{block}

\begin{block}{Block 3}
Suspendisse tincidunt sagittis gravida. Curabitur condimentum, enim sed venenatis rutrum, ipsum neque consectetur orci, sed blandit justo nisi ac lacus.
\end{block}
\end{frame}

%------------------------------------------------

\begin{frame}
\frametitle{Multiple Columns}
\begin{columns}[c] % The "c" option specifies centered vertical alignment while the "t" option is used for top vertical alignment

\column{.45\textwidth} % Left column and width
\textbf{Heading}
\begin{enumerate}
\item Statement
\item Explanation
\item Example
\end{enumerate}

\column{.5\textwidth} % Right column and width
Lorem ipsum dolor sit amet, consectetur adipiscing elit. Integer lectus nisl, ultricies in feugiat rutrum, porttitor sit amet augue. Aliquam ut tortor mauris. Sed volutpat ante purus, quis accumsan dolor.

\end{columns}
\end{frame}

%------------------------------------------------
\section{Second Section}
%------------------------------------------------

\begin{frame}
\frametitle{Table}
\begin{table}
\begin{tabular}{l l l}
\toprule
\textbf{Treatments} & \textbf{Response 1} & \textbf{Response 2}\\
\midrule
Treatment 1 & 0.0003262 & 0.562 \\
Treatment 2 & 0.0015681 & 0.910 \\
Treatment 3 & 0.0009271 & 0.296 \\
\bottomrule
\end{tabular}
\caption{Table caption}
\end{table}
\end{frame}

%------------------------------------------------

\begin{frame}
\frametitle{Theorem}
\begin{theorem}[Mass--energy equivalence]
$E = mc^2$
\end{theorem}
\end{frame}

%------------------------------------------------

\begin{frame}[fragile] % Need to use the fragile option when verbatim is used in the slide
\frametitle{Verbatim}
\begin{example}[Theorem Slide Code]
\begin{verbatim}
\begin{frame}
\frametitle{Theorem}
\begin{theorem}[Mass--energy equivalence]
$E = mc^2$
\end{theorem}
\end{frame}\end{verbatim}
\end{example}
\end{frame}

%------------------------------------------------

\begin{frame}
\frametitle{Figure}
Code to include your own image as figure~\ref{fig:acis} from the `Figure' directory.
\begin{figure}
\includegraphics[width=0.5\linewidth]{acis}
\caption{Acis logo}
\label{fig:acis}
\end{figure}
\end{frame}

%------------------------------------------------

\begin{frame}[fragile] % Need to use the fragile option when verbatim is used in the slide
\frametitle{Citation}
An example of the \verb|\cite| command to cite within the presentation:\\~

This statement requires citation \cite{Weng98}.
\end{frame}

% Final frame, subtext is optional
% Note: should be plain
%\setbeamertemplate{final page}[rwth][Any questions?]{Thank you for your attention}
%\begin{frame}[plain]
%\usebeamertemplate{final page}
%\end{frame}

\section*{References}
\begin{frame}[allowframebreaks]
  \Large{
    \bibliographystyle{apalike}
    \bibliography{../../bibliography/bibliography}
  }
  \end{frame}

\end{document}
